\documentclass[12pt]{article}
\usepackage[utf8]{inputenc}
\usepackage[T1]{fontenc}
\usepackage[margin=1in]{geometry}
\usepackage{setspace}
\usepackage{xcolor}
\usepackage[colorlinks=true, urlcolor=blue!60!black, linkcolor=blue!60!black]{hyperref}
\setlength{\parindent}{2em}
\setlength{\parskip}{1.25em}
\renewcommand{\baselinestretch}{1.0}
\begin{document}
\begin{center}
{\LARGE \textit{Haaland v. Brackeen}, 599 U.S. 255 (2023)\par}
\vspace{0.8em}
{\large Case Brief --- Four-Page Edition\par}
\end{center}
\vspace{1em}

\section{Memory Jogger}

Rejecting a multi-front attack on the Indian Child Welfare Act, the Court held 7--2 that ICWA falls within Congress's Article~I power over Indian affairs, that its custody rules do not commandeer the States, and that the individual challengers and Texas lacked standing to press equal protection claims against federal officials (\href{https://supreme.justia.com/cases/federal/us/599/21-376/}{599 U.S. 255}).

\section{Detailed Case Facts}

Congress enacted ICWA in 1978 after finding that ``an alarmingly high percentage of Indian families are broken up by the removal, often unwarranted, of their children from them by nontribal public and private agencies.'' In state custody proceedings involving an ``Indian child,'' ``any party'' seeking foster placement or termination must give notice, prove ``active efforts'' to prevent the family's breakup, and offer expert testimony of likely ``serious emotional or physical damage''; tribes may intervene; adoptive placements prefer extended family, tribal members, and then ``other Indian families'' absent ``good cause''; and States must keep placement records.

The suit arose from three custody disputes: the Brackeens' adoption of A.L.M. in Texas and later effort to adopt his sister; Altagracia Hernandez's voluntary placement of Baby O. with the Librettis; and the Cliffords' failed effort to adopt Child P. in Minnesota after a tribe enrolled her and she was placed with her grandmother.

\section{Procedural History}

The individuals, joined by Texas, Indiana, and Louisiana, sued federal officials in the Northern District of Texas; tribes intervened. The district court ruled for plaintiffs on equal protection, Tenth Amendment, nondelegation, and APA grounds. A Fifth Circuit panel reversed; the fractured en banc court upheld Congress's power but struck several provisions under the Tenth Amendment and divided evenly on others. The Supreme Court heard argument on November 9, 2022, and decided the case on June 15, 2023.

\section{Judicial Votes}

Affirmed in part, reversed in part, and vacated in part, \textbf{7--2}.
\begin{itemize}
\item \textbf{Majority:} Justice Barrett, joined by Chief Justice Roberts and Justices Sotomayor, Kagan, Gorsuch, Kavanaugh, and Jackson.
\item \textbf{Concurrences:} Justice Gorsuch (joined in part by Sotomayor and Jackson); Justice Kavanaugh.
\item \textbf{Dissents:} Justice Thomas; Justice Alito.
\end{itemize}

\section{Holding}

(1) ICWA is within Congress's Article~I authority over Indian affairs; there is no ``family law carveout.'' (2) No challenged provision commandeers the States: \S~1912 applies evenhandedly, \S~1915 imposes no search duty on state agencies and validly preempts state law in state courts, and the recordkeeping duties are permissible ``ancillary'' judicial tasks. (3) Neither the individuals nor Texas had standing to raise the equal protection or nondelegation claims.

\section{Analysis of Opinions}

\textbf{Majority (Barrett).} Congress's ``plenary and exclusive'' power rests on the Indian Commerce Clause, the Treaty Clause, constitutional structure, and the trust relationship. Valid federal law may preempt state family law, and petitioners never asked to overrule the precedents that foreclosed their narrower arguments. Under \textit{Murphy} and \textit{Reno v. Condon}, \S~1912 binds ``any party,'' including private litigants; under \textit{Adoptive Couple}, \S~1915 makes the objecting party, not the State, find alternatives; and early statutes confirmed in \textit{Printz} allow ancillary court recordkeeping. The individuals' injury was not redressable, because state officials who apply the preferences are not bound by a judgment against federal officers. Texas has no equal protection rights of its own and cannot sue the United States as \textit{parens patriae}.

\textbf{Gorsuch (concurring).} Recounting boarding schools and state removals, he described an originalist ``Indian-law bargain'' in which tribes keep sovereignty and States have no authority over Indian affairs.

\textbf{Kavanaugh (concurring).} The equal protection issue remains open and ``is serious.''

\textbf{Thomas and Alito (dissenting).} Thomas argued that ``Commerce'' is economic in all three clauses and that plenary power is ``judicial ipse dixit''; Alito contended that federal power cannot subordinate children's best interests in state family courts.

\textbf{Subsequent developments.} Days later, \textit{United States v. Texas} (2023) quoted the rule that ``it is the judgment, not the opinion, that demonstrates redressability,'' and \textit{Murthy v. Missouri} (2024) cited \textit{Brackeen} to confirm that States lack \textit{parens patriae} standing against the Federal Government.

\section{Examples: Future Applications}

\begin{enumerate}
\item \textbf{Same-Side (Upheld): Tribal Members' Estates.} A federal rule requiring ``any party'' in state probate to notify the tribe and apply a tribal-relative preference fits Article~I and evenhandedness.
\item \textbf{Same-Side (No Standing): Agency Suit.} An adoption agency suing the Secretary over \S~1915(b) lacks redressability; its families may sue in state court.
\item \textbf{Opposite-Side (Invalid): Nationwide Mandate.} A child-welfare law for all children, resting on the Interstate Commerce Clause and commanding state agencies alone, fails \textit{Lopez}, \textit{Printz}, and \textit{Murphy}.
\item \textbf{Opposite-Side (Likely Invalid): State-Only Search Duty.} A search-and-report duty addressed only to state agencies is a direct command to state officers.
\item \textbf{Fence-Sitter: Equal Protection in State Court.} Foster parents challenge the ``other Indian families'' preference under \textit{Palmore}; the tribe invokes \textit{Morton v. Mancari}.
\end{enumerate}

\section{Critique}

\textbf{Unstated limits and formalism.} The majority conceded that the precedent is ``unwieldy'' yet declined to define the borders of Congress's sphere, a contrast with \textit{Lopez}. Critics warn that letting Congress bypass \textit{Printz} by addressing ``any party'' elevates phrasing over reality where state agencies bring most proceedings.

\textbf{Competing originalisms.} Justices Gorsuch and Thomas both relied on Founding-era dictionaries and statutes yet disagreed on whether ``Commerce with the Indian Tribes'' was a diplomatic term of art or an economic one.

\section{Key Quotations}

\begin{quote}
``It is a federal court's judgment, not its opinion, that remedies an injury; thus it is the judgment, not the opinion, that demonstrates redressability. The individual petitioners can hope for nothing more than an opinion, so they cannot satisfy Article III.''
\end{quote}
\noindent --- Justice Barrett, Opinion of the Court, 599 U.S. 255.

\begin{quote}
``It could not be otherwise---Article I gives Congress a series of enumerated powers, not a series of blank checks. Thus, we reiterate that Congress's authority to legislate with respect to Indians is not unbounded. It is plenary within its sphere, but even a sizeable sphere has borders. \ldots{} But the Constitution does not erect a firewall around family law.''
\end{quote}
\noindent --- Justice Barrett, Opinion of the Court, 599 U.S. 255.

\begin{quote}
``To be clear, however, `plenary' does not mean `free-floating.' A power unmoored from the Constitution would lack both justification and limits. So like the rest of its legislative powers, Congress's authority to regulate Indians must derive from the Constitution, not the atmosphere.''
\end{quote}
\noindent --- Justice Barrett, Opinion of the Court, 599 U.S. 255.

\end{document}
